\begin{itemize}
\item \textbf{Taxonomy and labels.} Class validity holds for the classes and labels of the calibration set. A class absent from the taxonomy, or labelled unreliably, is not protected. ECG5000 labels come from automated annotation, so on that benchmark ``class'' is a noisy proxy.
\item \textbf{Exchangeability.} We assume that objects of a class are exchangeable between calibration and test. On ZTF and PLAsTiCC this mirrors FLARE's own protocol. Spectroscopically labelled calibration objects can be selected on brightness, which correlates with the forecast target. If test objects are selected differently, coverage can drift, for every method. A survey whose class mix or selection changes over time needs recalibration. Class-conditional validity is invariant to changes of class priors alone \citep{podkopaev2021distribution}.
\item \textbf{Finite samples.} The default \gamut{} fits its weights, residual laws and temperature on the calibration objects and has no finite-sample guarantee; the uniform-convergence bound in \cref{app:proofs} is too loose to help at our class sizes. \gamut-cert is valid in finite samples under a strict split, but it is conservative when a class has few objects. A class with fewer than $1/\alpha-1$ certification objects in a horizon cell gets no guarantee there and is left out of that cell's threshold; at 95\% on small calibration sets these are exactly the rare classes.
\item \textbf{Horizon cells.} With few calibration objects, the merging rule coarsens the horizon cells, at the extreme to a single cell. Validity is then per merged cell, not per FLARE bin. The realised cells are listed in \cref{app:cert}, and per-bin coverage is checked by bootstrap.
\item \textbf{Empty bands.} A level set can be empty at an input whose class hypotheses are all diffuse; the length-optimal band then gives up that point. On ZTF this happens for at most 0.01\% of future points. Elsewhere it is not negligible: at 68\% up to 7.8\% of points on PLAsTiCC-train, 4.6\% on ElectricDevices and 1.4\% on PLAsTiCC-DDF, and at 95\% up to 1.5\% on PLAsTiCC-train and 1.2\% on ECG5000 (\cref{app:unfavourable}). We count empty bands as misses. A deployment that needs a band at every point can fall back to the posterior-mixture band there, which we did not evaluate.
\item \textbf{Discrete targets.} On ElectricDevices 52--59\% of the targets are exactly zero, so $Y$ has no Lebesgue density there, against the assumption of \cref{prop:oracle}, and band lengths depend on the evaluation grid near zero.
\item \textbf{Length on a grid.} Coverage is computed from the exact score, and length on a grid of 400 uniform points plus 49 points around every class hypothesis. A band narrower than the finest grid cell is counted as empty, and hence as a miss, which is conservative.
\item \textbf{Unions of intervals.} When class forecasts disagree, the band can be a union of intervals. This is informative (``either it fades like a SN~Ia or stays bright like an SLSN''). In our benchmarks the hull is at most 2\% longer than the union on average, so the bands are almost always single intervals.
\item \textbf{Shared base models.} We changed only the calibration step. A better base forecaster or classifier narrows every method's bands. For compute reasons we used a learning rate of 0.05 (ZTF) or 0.1 (other benchmarks) instead of FLARE's 0.02, the same for all methods.
\item \textbf{UCR normalisation.} The archive z-normalises each series over its full length, which leaks scale information into the prefix. All methods see the same inputs, so their comparison is unaffected, but absolute widths on UCR are optimistic.
\item \textbf{Test-set exposure and the validity rule.} Several design choices (adaptive cells, the support of the residual laws, the evaluation grid, the validity rule) were revised after test results had been seen (\cref{app:history}). The pre-registered held-out set (\cref{app:held}) bounds this effect, but it is one benchmark under the cross-fitted protocol, with six settings. The bootstrap rule tests each class separately, without a multiplicity correction, and it has little power for classes with few objects: a class of 20 objects is flagged only when its coverage is far below nominal. Both points make the reported validity rates optimistic, for \gamut{} and for every baseline alike.
\item \textbf{Same-family review.} The internal reviews of this work were done by models of the same family as the executor (\texttt{review\_independence: same-family}); they are not independent peer review.
\end{itemize}
